% Generated draft: requires mathematical and source review.
\documentclass[11pt]{article}
\usepackage[margin=25mm]{geometry}
\usepackage{fontspec}
\setmainfont{Latin Modern Roman}
\usepackage{amsmath,amssymb,mathtools,unicode-math}
\setmathfont{Latin Modern Math}
\usepackage{enumitem,needspace,xurl,hyperref}
\setlength{\parindent}{0pt}
\setlength{\parskip}{5pt}
\setlength{\emergencystretch}{3em}
\newcommand{\N}{\mathbb{N}}
\newcommand{\Z}{\mathbb{Z}}
\newcommand{\Q}{\mathbb{Q}}
\newcommand{\R}{\mathbb{R}}
\newcommand{\C}{\mathbb{C}}
\newcommand{\F}{\mathbb{F}}
\newcommand{\A}{\mathbb{A}}
\newcommand{\PP}{\mathbb{P}}
\newcommand{\E}{\mathbb{E}}
\newcommand{\eps}{\varepsilon}
\newcommand{\dd}{\,\mathrm{d}}
\begin{document}
\section*{UnsolvedMath ID: 800313}
\textbf{Source problem code:} catalogue entry
\subsection{Definitions and mathematical statement}
Let \(S^1=\R/\Z\), and write \(R_\alpha(x)=x+\alpha\pmod{\Z}\). For an orientation-preserving homeomorphism \(f:S^1\to S^1\), choose a lift \(F:\R\to\R\) satisfying \(F(x+1)=F(x)+1\). Its rotation number is
\[
\rho(f)=\left[\lim_{n\to\infty}\frac{F^n(x)-x}{n}\right]\in\R/\Z,
\]
independent of \(x\) and of the chosen lift modulo \(\Z\). The rotation number is irrational if it has a representative \(\alpha\in\R\setminus\Q\).

The usual irrational-rotation conclusion of Denjoy's theorem is that there exists a homeomorphism \(h:S^1\to S^1\) such that
\[
h\circ f\circ h^{-1}=R_\alpha,
\qquad [\alpha]=\rho(f).
\]
Using this interpretation of the reference to Denjoy's theorem, the supplied question asks: does every orientation-preserving circle homeomorphism \(f\) with irrational rotation number and real-analytic graph
\[
\Gamma_f=\{(x,f(x)):x\in S^1\}\subset S^1\times S^1
\]
satisfy this conjugacy conclusion?

The source does not explicitly state the orientation and irrational-rotation hypotheses; these are made explicit here to identify the usual Denjoy conclusion. It also does not define ``real-analytic graph.'' In particular, it does not specify whether \(\Gamma_f\) must be a real-analytic embedded one-dimensional submanifold of the standard real-analytic torus, or may be a real-analytic subset with singularities. Neither interpretation is silently selected, and real analyticity of the map \(f\) is not substituted for the stated graph condition. The intended graph class therefore requires clarification before this can be regarded as a fully specified mathematical statement.
\subsection{Short English statement}
Navas asks whether a circle homeomorphism with a real-analytic graph satisfies Denjoy's conclusion: under the usual orientation-preserving and irrational-rotation hypotheses, it is topologically conjugate to the corresponding rigid rotation. The source leaves the precise meaning of a real-analytic graph unspecified.
\subsection{Sources}
\begin{itemize}
\item Supplied local UnsolvedMath registry record: permanent ID \(800313\), source problem code \texttt{catalogue entry}, title \emph{Group Actions on 1-Manifolds --- Question 30}. The record attributes the question to Navas (2017) and labels it \texttt{open}; this is a supplied status label, not an independently verified assessment.
\item Navas, \emph{Group Actions on 1-Manifolds: A List of Very Concrete Open Questions} (2017), \url{https://arxiv.org/abs/1712.06462}. The supplied provenance identifies the item as ``Question 30 (\#29),'' records access on 2026-10-16, and gives the extraction method as \texttt{arxiv\_latex}. These provenance details are supplied by the registry and were not independently verified here.
\end{itemize}
\end{document}
